\section{Worked Examples and Probe Evidence}

\tool is evaluated through interface coverage and diagnostic separation. The
evidence has two layers. First, a same-item Musical
contrast compares GRID-style RQ-KMeans and ReSID/GAOQ under one adapter
contract, which is the right setting for the addressability-versus-alignment
finding. Second, LETTER and LC-Rec exercise released item-index artifacts,
showing that the same contract is not tied to the Musical case study. Controls
and mechanism probes then separate structural, behavioral, and capacity-driven
signals.

\paragraph{Worked-example protocol.}
The case study uses the same 23,742 Musical items for the GRID and ReSID rows.
GRID is a controlled feature-text export: ReSID processed feature text is
embedded and passed through a released GRID/RQ-KMeans-style residual k-means
module. ReSID is a bounded FAMAE-to-GAOQ export on the same item universe. The
protocol supports a same-item diagnostic contrast; end-to-end pipeline
reproduction is outside this artifact-layer protocol.
In Table~\ref{tab:musical-diagnostic}, D2 is full-code aliasing, D3 is level-1
weighted co-occurrence-prefix recall, D4 is tail unique-\sid ratio, and D5
lists active prefix counts by depth. Daggered rows are structurally
item-unique; bold rows mark the primary same-item contrast and italics mark
controls.

\begin{table}[t]
\caption{Musical diagnostic profile.}
\label{tab:musical-diagnostic}
\scriptsize
\setlength{\tabcolsep}{1.8pt}
\begin{tabular*}{\columnwidth}{@{\extracolsep{\fill}}lrrrrl@{}}
\toprule
Artifact & Unique SIDs & D2 & D3 & D4 & D5 levels \\
\midrule
\multicolumn{6}{@{}l}{\emph{RQ-style exports}} \\
\textbf{GRID-style ft} & 3,749 & 0.977 & 0.055 & 0.370 & 64/3.4k/3.7k \\
GRID-style cap & 9,874 & 0.779 & 0.080 & 0.639 & 32/9.3k/9.9k \\
RQ-min ref & 17,247 & 0.440 & 0.065 & 0.883 & 32/2.4k/17.2k \\
\addlinespace[1pt]
\multicolumn{6}{@{}l}{\emph{ReSID/GAOQ export}} \\
\textbf{ReSID}\textsuperscript{\dag} & 23,742 & 0.000 & 0.154 & 1.000 & 32/1.3k/23.7k \\
\addlinespace[1pt]
\multicolumn{6}{@{}l}{\emph{Controls}} \\
\emph{Cat-prefix}\textsuperscript{\dag} & 23,742 & 0.000 & 0.447 & 1.000 & 30/83/313/23.7k \\
\emph{Pop-balanced} & 22,707 & 0.086 & 0.303 & 0.962 & 4/1.0k/22.7k/22.7k \\
\emph{Hash-collide} & 256 & 1.000 & 0.004 & 0.032 & 256/256/256/256 \\
\bottomrule
\end{tabular*}
\end{table}

\begin{table}[t]
\caption{Mechanism-probe calibration.}
\label{tab:mechanism-probes}
\scriptsize
\setlength{\tabcolsep}{2pt}
\begin{tabularx}{\columnwidth}{@{}>{\raggedright\arraybackslash}p{0.29\columnwidth}>{\centering\arraybackslash}p{0.11\columnwidth}>{\raggedright\arraybackslash}X>{\raggedright\arraybackslash}X@{}}
\toprule
Mechanism & D & Baseline signal & Activated signal \\
\midrule
Qualified aliasing & D2 & hash 1.19x & co-occur 3.86x \\
Capacity budget & D1/D4 & head 1.000 & tail 0.028 \\
Variable depth & D5 & max-depth 12,010 & active 7,914 \\
\bottomrule
\end{tabularx}
\end{table}

Table~\ref{tab:musical-diagnostic} makes the same-item rows comparable under
one schema. The GRID and ReSID rows are the tokenizer contrast; the capacity
ablation, RQ-min reference, and controls explain which parts of the contrast
come from capacity, addressability, or prefix semantics. In this setup, GRID-style ft
exposes high aliasing pressure and limited tail capacity, while the ReSID/GAOQ
row exports one full code per item and has higher D3-L1 collaborative prefix
recovery. The daggered rows mark a structural floor: collision-free full codes
capture addressability, while recommendation quality depends on downstream
generation and ranking.

The capacity-matched GRID row tests whether the original GRID pressure is only
a budget artifact. With a larger 32/1280/1280 budget, unique full codes rise to
9,874 and D2 aliasing drops to 0.779, but the row remains far from the
structurally item-unique ReSID and category-prefix rows. Capacity explains part
of the pressure, while the non-item-unique leaf keeps the contrast at the
artifact-profile layer. RQ-min is a
minimal residual-quantization reference adapter that passes the same validator;
it tests adapter independence on the full Musical item universe and is reported
as a reference row rather than an additional tokenizer line.

\paragraph{Diagnostic findings.}
The worked example supports one central finding: addressability is not
behavioral prefix alignment. ReSID is structurally item-unique, and the
category-prefix control is also alias-free, but D3 separates their profiles.
ReSID is a learned export, whereas the category row is a deterministic control
for interpreting the metric scale.
The non-learned category-prefix row recovers more level-1 co-occurrence
neighbors (0.447) than ReSID (0.154) or either GRID row (0.055--0.080). The
result locates a prefix-alignment gap: learned or item-unique addresses can
still fail to group behaviorally related items. This failure mode is visible
only when the mapping is inspected directly.

The supporting evidence is diagnostic separation. Across the RQ-style rows,
D2 aliasing falls from 0.977 (GRID-style ft) to 0.779 (GRID-style cap) and 0.440
(RQ-min), while D3 stays in the narrow 0.055--0.080 range. Capacity and
aliasing can improve without producing behaviorally aligned prefixes, making
D2 and D3 complementary diagnostics. The ablation weakens a simple ``too little
capacity'' objection without making the row item-unique. Table~\ref{tab:mechanism-probes}
gives the same point as controlled checks: qualified aliasing separates
co-occurrence aliases from hash aliases, capacity probes separate head and tail
allocation, and variable-depth probes separate maximum-depth from realized
prefix cost.

Released item-index adapters test the resource beyond the same-item case.
LETTER and LC-Rec Instruments artifacts both pass the same validation path on
9,922 items; both expose 9,897 unique full \sid{}s, with D3-L1 0.109 for LETTER
and 0.052 for LC-Rec. They appear in the catalog rather than the same-item
table because they use a different item universe; their role is to show that
\tool ingests released tokenizer mappings as well as locally generated
contrasts.

Additional checks bound the D3 claim and are indexed in the released
reproducibility matrix. On an All\_Beauty 20k vertical, a coarse
category-prefix control reaches D3-L1 0.968 while GRID/RQ-KMeans feature-text
rows are 0.081, 0.087, and 0.090 across seeds 42--44. Here D3 measures
taxonomy--co-occurrence alignment under the available coarse category signal,
which differs across datasets. The high All\_Beauty value therefore remains a diagnostic signal:
the available coarse taxonomy is more aligned with observed co-occurrence there
than in Musical. Two unified 5,000-user fixed-reranker probes sharpen the
boundary. Across eight Musical rows and six All\_Beauty rows, D3 is strongly
associated with candidate target recovery under train-only prefix retrieval
(depth-1 Spearman 0.976 and 1.000). At \(K=20\), the corresponding Spearman
values with ranked Recall/NDCG are 0.970/0.952 on Musical and 0.657/0.657 on
All\_Beauty. D3 exposes candidate-exposure structure before the reranker or
generator decides final ordering. A Sports 20k GRID export has
complete metadata and interaction joins, D3-L1 0.055, duplicate-\sid rate
0.592, and D5 prefixes 128/7986/8165. Sports, MovieLens, DACT, LETTER, and
LC-Rec rows exercise portability and adapter coverage around the same-item
Musical finding.
